\ifdefined\MyWebBook\else
\documentclass[../textbook.tex]{subfiles}
\begin{document}
\fi
\bookchapter{模型安全}{ch:model-application-security}

模型应用安全涵盖模型输入输出、数据访问、身份认证与动作执行。工具调用需要结构校验和执行授权，敏感信息保护需要覆盖检索、上下文、缓存与日志。威胁分析以这些组件构成的完整系统为对象。

本章承接\bookxref{ch:rag-systems}中的知识与权限边界，把相同原则扩展到训练数据、模型制品、工具执行和运行治理。威胁模型描述攻击路径，控制机制限制可达后果，运行机制负责监测与恢复。

\section{安全问题的对象与范围}

\subsection{可靠性、行为安全、防护和隐私}
\term{可靠性}{Reliability}{}关注系统在给定条件下能否持续完成预期功能。\term{行为安全}{Safety}{}关注失败和误用是否造成不可接受的伤害。\term{安全防护}{Security}{}关注主动对手能否破坏资产、权限和系统行为。\term{隐私}{Privacy}{}关注个人信息如何被收集、使用、保留和披露。这些目标相互关联，却没有共同的单一分数。

例如，系统可能稳定执行错误授权，这是一种可靠的越权行为；模型也可能在无人攻击时给出有害建议。加密传输保护的是链路，已授权服务把过量数据放入提示属于另一层问题。准确率、内容拒绝率、权限控制和数据最小化分别回答不同问题。

\term{保障论证}{Assurance Case}{}将目标声明、适用假设、控制和证据连接起来。例如“租户A不能读取租户B的文档”应由服务身份绑定、数据层授权、缓存隔离及对应测试支持；系统提示里的一句话提供不了这一保证。论证同时记录残余风险、责任人和失效时的处理方式，使结论有明确边界。

本章以 NIST AI RMF 1.0（2023年）与生成式人工智能配置文件 NIST AI 600-1（2024年）作为风险治理参考，以 OWASP 2025版相关条目辅助检查应用风险\cite{nist2023airmf,nist2024genaiprofile,owasp2025promptinjection}。风险分析需结合具体资产、部署场景与适用要求开展。

\subsection{威胁影响}
\term{威胁模型}{Threat Model}{}说明要保护的资产、对手能力、攻击入口、信任边界及可能后果。资产不仅是权重和用户数据，还包括分词器、模板、适配器、索引、会话记忆、工具凭据、日志和算力预算。攻击者可能只控制一页公开网页，也可能能提交训练样本、修改软件依赖或使用一个普通租户账号；不同能力会产生不同可达路径。

\term{信任边界}{Trust Boundary}{}是信息或能力从一个授权与可信假设范围进入另一个范围的位置。外部文件进入解析器、检索文本进入模型、模型建议进入执行器，都是需要独立检查的边界。各处理阶段按自身用途重新检查授权。允许读取某份文件与允许将其全文发送给第三方是两项不同的授权。

治理先确定使用边界、责任和风险接受权限；场景分析识别资产与后果；测量通过评测、红队和运行数据形成证据；管理则决定缓解、限制能力、接受残余风险或停止使用。四类活动相互反馈，每次变化都要重走其中的相关环节。每次数据、模型或工具变化都应明确哪些既有证据仍有效，哪些需要补充。

风险登记应先描述不可接受后果，再沿系统数据流寻找可达路径。一个完整条目包含资产、主体、前提、入口、影响、现有控制、验证证据、负责人和处理状态。空白条目表示尚未分析，风险本身依然存在。扩展模型能力、接入工具、更换检索源和改变用户群都会改变威胁模型，每次变化后都需重新登记。

\begin{assumption}[本章风险与授权模型的前提]
概率推导以明确的请求总体、时间窗口和损失单位为前提；授权由可信身份和资源状态提供，模型自报的身份字段不参与判定，也无权提升自身权限。攻击者可控制其获准提交的输入和部分外部内容，但不能任意修改受保护的授权引擎；若该假设失效，必须另行分析控制平面的失陷。
\end{assumption}

\begin{table}[htbp]
\centering
\caption{模型应用安全与风险治理核心数学符号}\label{tab:rev-security-notations}
\begin{tabularx}{\linewidth}{p{24mm}X}
\toprule 符号 & 含义与范围\\\midrule
$u,a,r,c$ & 已认证主体、动作、资源与运行上下文。\\
$P_v,\operatorname{Allow}$ & 版本为 $v$ 的策略及其资源级授权判定。\\
$\ell,L_i,p_i$ & 随机损失、场景损失和概率，$L_i\ge0$、$p_i\in[0,1]$。\\
$F_j$ & 第 $j$ 个控制层未阻止指定攻击的事件。\\
$w_i,\mathcal A_i$ & 预先确定的风险权重和第 $i$ 个样本的允许决策集合。\\
$N,x,\alpha$ & 评测试验数、失败数和单侧置信界显著性水平。\\
$h_a,k$ & 最终动作摘要与幂等键，分别绑定动作内容和逻辑执行请求。\\
\bottomrule
\end{tabularx}
\end{table}

\section{风险度量边界}

\subsection{期望损失与主观等级}


若场景构成互斥划分，场景 $i$ 发生概率为 $p_i$，其条件平均损失为 $L_i$，则
\begin{equation}
 \mathbb E[\ell]=\sum_i p_iL_i.
 \label{eq:security-risk}
\end{equation}
若各类损失是可相加的不同分量，也可利用期望的线性性分别求和；同一后果被多个攻击路径重复计数会高估期望。损失单位可以是同一种货币或其他明确计量；生命、隐私与业务影响难以合理折算，应分别呈现。

\begin{example}[风险估计与不确定性]
例如构造一个季度风险模型：某类事件概率0.02，发生后的平均恢复损失为50000个成本单位，则该分量的期望为1000。若概率只知道落在0.01至0.04，损失在30000至80000，则可给出300至3200的粗范围；只报1000的点估计会隐去这层不确定性。低概率而不可接受的后果还可设独立能力限制；组织对不可接受后果的判断另有标准，不进入平均损失。
\end{example}

“可能性1至5、严重性1至5、暴露程度1至5”的评分属于序数尺度。乘积可作为预先约定的排序启发式，但不具备真实期望损失语义，排序甚至对保持顺序的重新编码不稳定。两个条目为 $(1,5)$ 和 $(2,3)$，乘积分别5和6；若只把最高等级5重新编码为100，等级顺序未变，乘积却变为100和6。这个反例说明评分结果依赖量表约定，乘积差一倍推不出风险差一倍。

\subsection{控制效果与相关失败}
设 $R_0,R_1$ 为同一总体在无某控制和有该控制时的期望损失，且 $R_0>0$，定义有效率 $e=1-\frac{R_1}{R_0}$，才有
\begin{equation}
 R_1=R_0(1-e).
\end{equation}
它是根据已定义量得到的恒等关系；难点是获得可比较的反事实估计。把专家打分的“有效性80分”直接替换为 $e=0.8$ 没有统计依据。若控制使攻击转移到别处，或改变了用户总体，两个风险的可比性也随之丧失。

\term{纵深防御}{Defense in Depth}{}使用多个控制点限制攻击影响。若攻击成功需要所有层都漏放，则
\begin{equation}
 \Pr(\cap_{j=1}^mF_j)=\Pr(F_1)\prod_{j=2}^m
 \Pr(F_j\mid F_1,\ldots,F_{j-1}).
 \label{eq:security-defense}
\end{equation}
只有独立失败时才可将条件概率替换为边缘概率。两个各有0.1漏放率的过滤器，独立时联合为0.01；若它们在相同攻击上一起失败，联合仍为0.1。对两个事件，在只知道边缘概率时有
\begin{equation}
 \max(0,p_1+p_2-1)\le\Pr(F_1\cap F_2)\le\min(p_1,p_2).
\end{equation}
共享失效模式会限制多次判断的收益。防护组合可覆盖不同机制，例如身份校验、资源范围限制、隔离执行、内容判断与审计，并在完整链路上验证。表\ref{tab:rev-safety-defense-layers}总结了大模型应用纵深防御的四层体系划分，表\ref{tab:rev-security-threat-matrix}梳理了大模型系统最核心的安全威胁矩阵及其控制手段。

\begin{table}[htbp]
\centering
\caption{大语言模型应用纵深防御四层体系对比}\label{tab:rev-safety-defense-layers}
\begin{tabularx}{\linewidth}{p{20mm}p{22mm}p{38mm}X}
\toprule
防御层次 & 部署介入位置 & 核心检测与阻断机制 & 防护局限与潜在旁路风险 \\\midrule
输入检测层 & 前置 API 网关 & 提示注入分类器、关键词过滤、文本规范化 & 易被对抗编码混淆、同音代换或隐蔽多语言绕过 \\
模型对齐层 & 模型权重与参数 & 安全强化学习（RLHF/DPO）、负样本微调拒答 & 存在对抗提示后缀、越狱诱导与内在知识反转风险 \\
输出审核层 & 流式后处理网关 & 敏感信息扫描（PII 遮蔽）、合规二次评审模型 & 增加首词元生成延迟，难以甄别隐蔽逻辑漏洞 \\
运行沙箱层 & 工具执行环境 & 容器资源隔离、网络白名单、临时最小权限凭据 & 虽增加系统编排复杂度，但构成拦截非授权动作的核心底线 \\
\bottomrule
\end{tabularx}
\end{table}

\begin{table}[htbp]
\centering
\caption{大语言模型核心安全威胁矩阵与纵深防御措施}\label{tab:rev-security-threat-matrix}
\begin{tabularx}{\linewidth}{p{22mm}p{28mm}p{24mm}X}
\toprule
威胁类别 & 典型攻击机制 & 潜在危害后果 & 核心缓解与防御策略 \\\midrule
直接与间接提示注入 & 用户诱导或外部数据污染 & 劫持模型决策流，执行非授权动作 & 严格区分指令与非受信数据，强加独立策略校验 \\
敏感信息出域泄露 & 提示携带密钥或模型记忆复现 & 造成隐私泄露与业务数据外逃 & 上下文数据最小化，外部传输强制脱敏与审计 \\
第三方供应链风险 & 恶意微调权重或被投毒依赖 & 隐藏后门触发器，执行任意代码 & 权重哈希签名校验，限制未验证动态插件加载 \\
数据与模型投毒 & 训练或微调语料植入暗门 & 特定触发词下激发非预期违规行为 & 数据来源清洗审计，异常数据密度与聚类检测 \\
非受限工具调用 & 赋予智能体过大系统与网络权限 & 破坏内部数据库，恶意批量发信 & 最小权限原则，高危敏感操作强制人工二次确认 \\
\bottomrule
\end{tabularx}
\end{table}

\section{提示注入和指令边界}

\subsection{直接注入与间接注入}
\term{提示注入}{Prompt Injection}{}利用模型可能把不可信数据解释为有效指令的机制，诱导其偏离授权任务。\term{直接提示注入}{Direct Prompt Injection}{}来自用户提交的输入；\term{间接提示注入}{Indirect Prompt Injection}{}来自网页、文件、检索材料或工具返回值等外部来源。研究已经展示间接内容能够影响集成语言模型的应用\cite{greshake2023indirect}。图\ref{fig:indirect-injection-flow}展示了典型的间接提示注入攻击链路与防御介入点。本章关心的是数据如何越过权限边界，某一组固定攻击措辞只是其中一种表现。

\begin{figure}[htbp]
\centering
\includegraphics[width=0.96\linewidth]{\subfix{../../figures/theory/indirect-injection-flow.pdf}}
\caption{间接提示注入（Indirect Prompt Injection）端到端攻击链路与纵深防御边界。攻击者通过投毒外部非受信数据（网页、PDF、邮件等），在检索增强生成（RAG）或工具调用读入模型上下文时劫持控制流，诱导越权工具调用或数据外发。}\label{fig:indirect-injection-flow}
\end{figure}

命令式文字要按对象与指令分开看。用户可能要求分析一份含指令的文档，此时文档内容是分析对象。反过来，文字看起来像普通说明也可能试图改变输出目的。字符关键词、语言和排版都表达不了完整的授权关系。图像内文字、音频转写和工具结果同样可能携带这类内容，跨模态输入处在同一条边界上。

\term{越狱}{Jailbreak}{}通常指诱导模型突破其行为约束；它与提示注入可能重叠，但不是完全相同问题。内容政策、数据出域与执行权限具有不同的失效路径。评估必须分别记录模型行为偏离和实际权限或数据影响。

\subsection{内容可信度与执行权限分开}
系统可以通过角色分隔、来源标记和训练改善模型对指令层级的理解，但文本标记只表达意图，与密码学身份是两种凭据。外部文档对认证主体、租户、工具白名单和审批规则没有修改能力；模型生成的“已批准”也提供不了批准凭据。权限保存在模型可读文本之外的受保护状态中。

数据最小化在输出检查之前生效。如果模型上下文含有它不需要的秘密，任何过滤失误都会扩大泄露后果。更强的边界是不给模型不必要的秘密；工具凭据由执行器临时取得，写进系统提示则会长期驻留在上下文中。系统提示可能包含知识产权或内部规则，却承担不了密钥库或唯一安全屏障的作用。

\begin{figure}[htbp]
\centering
\begin{tikzpicture}[node distance=6mm,box/.style={draw,align=center,text width=29mm,minimum height=12mm,font=\small},>=stealth]
\node[box] (x) {外部内容\\网页、文件、工具结果};
\node[box,right=of x] (m) {模型与编排\\生成动作建议};
\node[box,right=of m] (p) {独立策略检查\\身份、范围、预算};
\node[draw,align=center,text width=18mm,right=7mm of p,font=\small] (stop) {返回状态\\停止执行};
\node[box,below=of p] (e) {受控执行器\\短期最小权限};
\node[box,left=of e] (d) {资源与结果\\数据层授权};
\node[box,left=of d] (a) {最小化审计\\版本、决策、摘要};
\draw[->] (x)--(m);\draw[->] (m)--(p);\draw[->] (p)--node[right,font=\small]{仅允许}(e);\draw[->] (p)--node[above,font=\scriptsize]{拒绝或待批准}(stop);\draw[->] (e)--(d);\draw[->] (d)--(a);
\draw[dashed] (5.1,-.8)--(5.1,.9) node[above,font=\small]{执行信任边界};
\draw[->] (d.north)--(m.south) node[midway,left,font=\small]{结果仍是数据};
\end{tikzpicture}
\caption{模型建议到执行能力之间的独立边界。身份和凭据由控制层提供，外部内容不能授予权限。}
\label{fig:security-boundaries}
\end{figure}

\subsection{风险传播链}
设内部知识助手面临以下攻击：用户只要求总结某项目的内部报告，检索系统找到一份附有恶意说明的文档。模型把说明当作额外任务，提出向外部地址发送完整报告。若编排层只检查调用是合法 JSON，并使用拥有广泛发送权限的服务账号执行，就会把文本误解转化为数据外发。

此失败链至少包含四个不同环节：外部内容影响任务目的；模型提出超出原请求的动作；策略层没有验证用途、收件范围和数据级别；执行凭据允许现实外发。关键词过滤覆盖第一步中的部分变体，后三步分别需要动作准入、目的地址限制和数据出域策略。即使模型仍提出错误建议，这些执行控制也可阻止该后果。

另一条路径不需要显式发送工具：模型生成包含外部资源的富文本，客户端自动加载时可能发起网络请求。因此输出消费位置也属于威胁模型，应限制自动远程资源加载、按目标格式进行安全渲染，并避免将敏感内容放入可被外部访问的地址参数。展示层同样会发起动作，只保护工具网关会留下这条路径。

\section{工具授权与执行事务}

\subsection{外部授权判定}
令 $u$ 为已认证主体，$a$ 为动作，$r$ 为规范化资源，$c$ 包含租户、时间、用途和环境。执行许可至少满足
\begin{equation}
 \operatorname{Execute}(u,a,r,c)
 \Longrightarrow \operatorname{Allow}_{P_v}(u,a,r,c)=1.
 \label{eq:security-allow}
\end{equation}
这是必要条件而非充分条件，还要检查参数模式、业务约束、预算与操作前置状态。主体来自认证会话，资源归属来自可信数据层，模型参数中的自报字段只是待核对的线索。用户授权、服务身份能力和当前任务范围应取交集；任选一个足够宽的权限会绕过另外两项约束。

\term{最小权限}{Least Privilege}{}要求只提供完成任务必要的资源与动作范围。工具接口宜以明确业务能力为边界；暴露任意命令执行器会把权限面扩大到整台主机。读取、修改、外发和权限变更具有不同影响，应使用独立能力和凭据范围。OWASP 2025版将过多功能、过宽权限和过度自治列入过度代理风险讨论\cite{owasp2025agency}。

\term{结构模式}{Schema}{}校验约束字段、类型和取值形式，但合法格式仍可能引用错误租户、过量资源或不允许的目的地。资源标识必须规范化后验证，并在执行时使用同一个规范结果。先校验一个地址、再跟随未经验证的重定向，或者先检查路径字符串、再解析到其他对象，都会破坏这一对应关系。

\subsection{动作级批准}
对需要批准的高影响或不可逆操作，应在内容、对象、参数和影响已经确定后展示给具有相应权限的人。已存在且覆盖当前动作的授权不需要重复索取；模糊任务目标的授权范围以原任务为界。批准记录可绑定规范化动作摘要
\begin{equation}
 h_a=H(a,r,\mathrm{params},u,\mathrm{scope},v,\mathrm{expiry}),
\end{equation}
其中 $H$ 表示确定的摘要构造。参数、资源或权限状态改变后，原批准可能不再有效，需要重新判定。摘要能绑定内容，批准本身仍须由可信身份签发和验证。

\term{检查与使用时差}{Time-of-Check to Time-of-Use}{TOCTOU}指检查状态与实际执行间状态发生变化的风险。例如文档所有权在批准后改变。执行器应在执行点重新验证权限和必要资源版本，或使用数据库条件写入、事务与能力凭证绑定约束；把过期检查当作当前事实正是这类漏洞的成因。

\subsection{幂等、超时和未知结果}
\term{幂等性}{Idempotency}{}要求同一逻辑请求的重复提交不重复产生业务副作用。它需要稳定幂等键 $k$、动作摘要和执行结果的绑定；客户端随机生成一个字符串只提供了键的形式。相同键对应不同参数应拒绝；同一请求重试应查询或复用原执行记录。

工具超时可能发生在副作用已经完成、结果尚未返回之后。“超时”与“未执行”是两种状态。状态应区分已准备、已授权、执行中、成功、明确失败和结果未知。未知状态先查对账或结果接口；用新幂等键重新发起会制造第二个副作用。补偿操作与真实回滚属于不同机制：已经发送的数据无法靠撤销调用收回，某些业务动作只能进行后续补偿。

\begin{algorithm}[受控工具执行]
\AlgoInput{已认证请求、模型动作建议和当前策略。}
\AlgoOutput{拒绝、待批准、执行结果或结果未知状态。}
\begin{enumerate}
 \item 从受保护会话取得主体与任务范围，解析模型建议；身份或租户的自报字段与会话状态冲突时，以会话状态为准。
 \item 规范化资源与参数，验证结构、范围、用途、资源归属及预算；策略明确拒绝或必需状态不可取得时停止。
 \item 计算最终动作摘要。若策略要求批准，验证批准主体、摘要、范围及有效期；没有满足条件的批准则返回待批准状态并停止。
 \item 以“租户与幂等键”的唯一约束，原子建立并认领绑定动作摘要的执行记录；只有认领成功的执行者进入发送步骤。已有记录时返回已知结果或查询执行状态；相同键不同摘要则拒绝，不通过新建动作绕过冲突。
 \item 在执行点重查当前授权与资源前置条件，以最小权限短期凭据调用指定工具，不向模型暴露凭据。
 \item 将成功、明确失败或结果未知写入账本。未知结果进入对账流程；重试遵守相同逻辑请求和预算，不重复创建副作用。
 \item 按数据披露策略处理工具结果，记录最小化审计证据，再返回编排层。工具返回文本不改变后续授权规则。
\end{enumerate}
该算法要求账本支持原子认领，远端工具对同一动作键提供幂等提交或可靠的结果查询。认领后崩溃和发送后超时均须通过原动作身份恢复；远端既不支持幂等、也不能核对结果时，应保持结果未知并停止自动重发；本地唯一约束管的是本地记录，宣称外部恰好执行一次还需要远端证据。持久执行、发件箱与隔离代数的机制见\bookxref{ch:agent-runtime}。

不变量是未经授权的动作不能到达执行器，批准与实际参数一致，并在上述远端契约成立时使同一逻辑请求的重试不产生额外副作用。系统中存在绕过网关的其他执行路径时，必须纳入同一控制。
\end{algorithm}

\section{数据与模型供应链}

\subsection{投毒、后门和评测污染}
\term{数据投毒}{Data Poisoning}{}通过改变训练或检索数据影响系统行为。\term{后门}{Backdoor}{}使模型在特定触发条件下产生攻击者希望的行为，同时可能在普通样本中表现正常。两者的关系取决于攻击路径：后门可以通过训练数据、参数或依赖引入，数据投毒也可能只降低整体质量而不形成触发机制。

数据清单应保存来源、取得依据、用途、时间范围、处理版本、敏感等级和内容标识，并建立样本到数据快照、训练任务与模型版本的关系。来源集中度、重复率、标签变化和分布异常可帮助发现问题，检测结果可用于定位可疑来源与样本，并结合溯源和复核处理。下载位置公开说明不了内容是否适合训练，也说明不了来源是否拥有授权。

\term{评测污染}{Evaluation Contamination}{}使测试材料进入训练或调参过程，从而削弱评测结论。安全数据尤其需要控制：反复使用同一攻击模板优化，可能只提高已知变体拦截率。应独立维护训练、开发、保留测试和历史回归集合；新事件一旦转为回归用例，它就属于已知攻击。

\subsection{制品性校验}
\term{制品完整性}{Artifact Integrity}{}关注实际加载内容是否与指定制品一致。内容哈希可以检测字节变化，但攻击者若同时替换文件和未受保护的摘要，哈希相符也没有来源保证。\term{数字签名}{Digital Signature}{}把制品声明绑定到签名主体，但仍需验证签名身份、信任链、签名用途和撤销状态。发布签名验证来源，行为风险通过模型与应用评估检查。

权重之外，配置、分词器、会话模板、适配器、处理器、自定义代码和运行依赖都会影响行为。\term{软件物料清单}{Software Bill of Materials}{SBOM}列出软件组件及版本关系，模型制品清单还应补充权重与数据、评测和处理器关系。OWASP 2025版供应链条目强调预训练模型、数据和第三方组件带来的风险\cite{owasp2025supplychain}。清单支持组件追溯，发布评估检查对应的行为与依赖。

Safetensors 采用以张量为中心的序列化方式，旨在避免依赖任意对象反序列化来加载权重\cite{huggingface2026safetensors}。它减少一类加载风险；同仓库的自定义代码、模板与模型中的行为后门仍需另行审查。远端代码执行能力应默认关闭；业务确实需要时，按固定代码版本审查，并在隔离构建与加载环境中限制凭据和网络。

\subsection{构建、评测与部署绑定}
不可变发布应把权重、配置、处理器、模板、依赖锁定文件和镜像摘要绑定为一个版本集合。评测清单引用的摘要必须与部署集合一致；“模型名相同”支撑不起复用证据。改变适配器或会话模板也可能改变安全行为，应按影响重新评估。

构建环境需要限制外部执行和秘密访问，加载环境需要限制文件与网络能力，运行环境需要限制服务身份。漏洞扫描只能检查其覆盖的已知问题，行为评测只能覆盖相应任务分布；二者应互补。回滚时必须明确是回滚权重、索引、模板、策略还是完整制品集合；只恢复旧权重而继续加载受污染的适配器，回滚就没有完成。

\section{隐私和多租户隔离}

\subsection{数据最小化}
\term{个人可识别信息}{Personally Identifiable Information}{PII}包括能够识别或关联个人的信息，其具体范围应按场景确定。隐私控制从明确用途和最小必要数据开始，再处理访问、保留和披露。正则脱敏可能漏掉间接标识，也可能删除任务所需内容；对一处输出脱敏改变不了整个数据生命周期的暴露面。

将姓名替换为稳定标识属于去标识处理的一种形式，但映射表或其他字段组合仍可能恢复关联。向量、摘要、缓存和调试日志也可能保留敏感信息；这些形式与公开数据之间没有自动通道。加密保护传输和存储状态，对获得解密权限的服务仍需数据范围和用途限制。

训练模型还可能记忆样本信息。删除原始记录可以阻止继续摄取，却不自动删除已训练参数中的影响；需要依据实际训练与删除机制分析是否重训、使用专门方法或限制模型使用。数据目录中的删除状态与参数层面的精确遗忘是两回事。

\subsection{多租户状态隔离}
租户标识来自服务端认证状态，并贯穿文档读取、向量索引、对象存储、会话记忆、答案缓存和遥测。向量检索后的最终过滤处在外部重排器之后，已经传出的未授权内容收不回来。缓存键还应考虑授权范围和策略版本；权限撤销应及时阻止旧缓存读取，普通到期时间在这里太迟。

会话记忆的写入同样是数据变更。外部文档的自我声明不构成长期系统规则；记忆应保留来源、作用域、有效时间和删除路径。跨会话共享摘要必须有明确的数据授权，模型判断“这条信息以后有用”在这里不是依据。

密钥应由专用系统向执行器提供短期、窄范围能力。生产、测试和不同工具的服务身份分开，可以限制凭据泄漏后的影响。审计日志默认记录标识、版本、决策与必要摘要，对敏感全文的临时取证另行控制访问和保留期。没有必要把完整提示和工具响应无差别复制到监控平台。

\subsection{训练记忆、成员推断与输出泄漏}
隐私泄漏需要按信息来源分别分析。推理时泄漏来自当前请求、检索结果、工具返回或历史记忆；训练时泄漏则可能来自权重对样本的记忆。Carlini等人在GPT-2上通过查询恢复了部分训练文本，包括个人可识别信息\cite{carlini2021trainingextraction}。该实验展示了 GPT-2 在相应查询条件下的训练文本泄漏路径。

\term{成员推断}{Membership Inference}{}判断某条记录是否参加过训练；训练数据提取尝试恢复记录内容。两者影响不同，例如仅暴露某人属于一个敏感群体的数据集，也可能造成隐私损害。验证必须分别报告成员判断的误报与漏报、原文恢复的匹配标准，以及攻击者预先掌握的信息；训练文本恢复需核对生成内容与实际训练记录的匹配。

防护应贯穿数据生命周期：入库前排除非必要秘密、核实用途并控制重复；训练与微调时限制敏感样本暴露，保留删除和重训路径；推理时先授权检索，再最小化进入上下文的字段；输出时按接收者权限检查答案、引用、附件与工具参数。\term{数据防泄漏}{Data Loss Prevention}{DLP}通过数据识别与披露策略限制敏感信息外发。规则匹配、实体识别和上下文分类可以互补，但须覆盖改写、跨轮拼接与间接标识；只检查身份证号等固定格式会漏掉这三类。

流式输出还存在不可收回的披露窗口：先发送全文再审核，无法撤销已经到达客户端的内容。高敏感字段应在释放前检查；需要跨句判断时可缓冲一个可判定的片段，必要时审核完整回答，并明确增加的延迟。检查范围包括可读思考摘要和中间工具结果。故障时应返回明确状态或限定回答；审核服务不可用时的默认行为不能是无限放行。

例如员工只能查询自己的报销记录，系统应由会话身份绑定员工编号，在数据层返回获准字段，随后计算金额。输出审核检查其他员工姓名与账户信息，日志只记录必要标识和结果摘要。仅在最终答案遮盖姓名保护不了已发送给外部模型、重排器或日志平台的完整记录。供应商是否用请求训练、保存多久、哪些功能保留状态，是三个需要分别核对的条件；私有部署同样需要内部访问和日志控制。

\subsection{差分隐私边界}
\term{差分隐私}{Differential Privacy}{DP}约束单个受保护单位对随机机制输出分布的影响。令相邻数据集 $D,D'$ 相差一个约定单位，$\mathcal M$ 为训练发布机制，$S$ 为任意可测输出集合，则 $(\varepsilon,\delta)$-差分隐私要求
\begin{equation}
 \Pr[\mathcal M(D)\in S]\le \conste{}^{\varepsilon}\Pr[\mathcal M(D')\in S]+\delta,
 \qquad \varepsilon\ge0,\quad 0\le\delta<1.
\end{equation}
相邻单位可以是记录，也可以是一个人的全部记录，二者保证不同。DP-SGD通过逐样本梯度裁剪、聚合时加噪声和累计隐私核算实现训练保护\cite{abadi2016dp}。仅做普通梯度裁剪或在答案中任意加噪声，达不到上述定义所要求的那种保证。

采用该方法需要报告相邻关系、采样方式、噪声与训练步数，以及最终隐私预算；额外的私有数据调参和发布也要纳入核算。它带来效用和计算成本取舍，作用范围是训练数据的统计隐私。推理时的文档授权与输出内容控制属于服务侧机制，需要分别验证。

\section{模型提取与未经授权的蒸馏}

\subsection{保护资产与可观察信息}
\term{模型提取}{Model Extraction}{}指通过访问目标模型等方式构建行为近似的替代模型。知识蒸馏本身是正常训练方法，其机制见\bookxref{ch:quantization}；本节讨论超出授权范围的数据采集和能力复制。应分别确定允许的调用、批量生成和训练用途；用户要求解释步骤本身说明不了滥用。

对手可能取得权重，也可能只能通过API观察文本。权重一旦交付到对手控制的环境，服务端限流和输出策略就约束不到其本地执行；开放权重的发布决策须接受这一控制边界。API服务可以隐藏权重、中间表示与完整词元分布，并控制采集规模，但只要持续提供有用答案，观察者就可能利用答案训练其他模型。防护目标应表述为限制暴露、提高采集成本和发现未经授权的行为；保证输出无法被学习不在可达目标之列。

设请求为 $x$，教师生成的推理文本为 $z$，答案为 $y$，学生参数为 $\phi$。在已取得且获准使用的数据上，答案监督和轨迹监督可分别写成
\begin{align}
 \mathcal L_{\mathrm{answer}}&=-\mathbb E_{(x,y)}\log p_\phi(y\mid x),\\
 \mathcal L_{\mathrm{trace}}&=-\mathbb E_{(x,z,y)}\log p_\phi(z,y\mid x).
\end{align}
期望针对相应训练样本分布；序列概率按词元自回归分解。第二个目标增加了中间步骤监督，但轨迹错误、冗长或与学生能力不匹配时也可能降低效果。隐藏 $z$ 可以减少直接可用的轨迹，仍留下答案监督和由答案另行构造解释的途径；这种事后解释与教师实际生成过的内部轨迹是两种东西。

\subsection{请求安全监测}
长期采集按数据用途与服务授权单独检查。运行控制可结合可信账号归属、组织级预算、并发与批量限制，以及跨请求的任务集中度、模板重复、调用节奏和账号关联。聚合分析必须有明确用途、最小必要字段和保留期限；共享出口地址或相似课程作业都可能形成正常聚集，单一关联信号支承不起定性结论。

Anthropic在2026年2月23日的公告中介绍了流量分类、行为指纹、跨账号协同检测和强化账号验证，并提到识别为构建训练数据而诱导思维链的行为\cite{anthropic2026distillationdefense}。公告描述了账号、流量与行为层面的检测方向，具体检测规则与效果应按实际系统评估。

工程上可按证据强度逐级限制额度、暂停批量能力或进入人工核查，并保留申诉和恢复路径。验收既要统计受控采集场景中的可获得样本量、检测延迟和漏检，也要检查正常研究、教学与企业批处理的误限制。用于行为检测的数据受其采集用途限制，复制到学生训练或其他无关用途会超出这一范围。故意向正常用户提供错误答案会损害产品可信度，不适合作为通用防护措施。

\section{思维链保护与思考签名}

\subsection{明文、摘要和受保护状态}
\term{思维链}{Chain of Thought}{CoT}是模型生成的中间推理文本，与神经网络全部内部计算的记录是两回事。安全设计应区分完整推理文本、面向用户的解释或摘要，以及供后续推理恢复的受保护状态。完整文本可能复述上下文中的敏感内容，也可能成为蒸馏监督；摘要减少直接披露，却仍需按数据权限审核。

用户需要核验结论时，可以提供依据、计算步骤、来源与不确定性；发布全部内部轨迹并不是核验的前提。解释质量应以可验证性评估，一段文字的语气连贯性与它是否忠实反映模型的因果过程无关。隐去原始轨迹也改变不了模型从最终答案泄露秘密的通道。

\subsection{Anthropic的signature字段}
按2026年9月11日核查的Claude Messages API官方Thinking文档，\texttt{thinking}块中的\texttt{signature}是承载加密完整思考内容的不透明字段，API在回传时用它验证思考块来自Claude。\texttt{display}设为\texttt{summarized}时返回可读摘要，设为\texttt{omitted}时可读字段为空，受保护内容仍用于续接；具体支持情况随模型而异\cite{anthropic2026thinking}。

集成时应按协议原样保存和回传需要的思考块，包括\texttt{redacted\_thinking}块；流式签名通过\texttt{signature\_delta}到达。加密字段按要求原样回传，解析为可读思维链或当作普通答案展示都会破坏其用途。该字段按服务方的状态回传协议使用，其密码算法与完整绑定范围以公开协议说明为准\cite{anthropic2026thinking}。

从控制作用看，加密限制无解密能力的观察者读取完整轨迹；来源与完整性验证限制伪造或篡改后冒充有效思考状态。签名验证来源与完整性，公开明文仍可被复制。即使完整CoT受保护，可读摘要、答案和工具轨迹仍可能提供监督；声称反蒸馏有效还需要相应采集与训练评估。

\begin{table}[htbp]
\centering
\caption{思考保护与反蒸馏控制的作用边界。}\label{tab:rev-security-thinking-protection}
\begin{tabularx}{\linewidth}{@{}lXX@{}}
\toprule 控制 & 直接作用 & 仍需单独处理的风险\\\midrule
签名或认证校验 & 核验来源与完整性 & 可读文本被复制、内容本身错误。\\
加密保存与回传 & 限制无密钥者读取状态 & 解密端权限、明文输出与日志泄漏。\\
摘要或省略展示 & 减少公开轨迹 & 摘要与答案的敏感信息及训练价值。\\
配额与行为检测 & 限制并识别规模化采集 & 误限制、分散采集及正常用途区分。\\
\bottomrule
\end{tabularx}
\end{table}

\subsection{应用侧的状态隔离与验收}
对于自建系统，若客户端只需恢复会话，可在服务端保存推理状态并返回受限句柄；需要客户端携带状态时，应使用经过审查的认证加密方案，同时在受保护上下文中绑定租户、会话、用途、版本与有效期，并设计重放处理。密文也应按敏感状态控制访问、日志与保留；不可读与可以跨租户共享是两回事。

验收应分为三个独立目标：用获准的合成敏感标记检查答案、摘要、工具参数与日志的泄漏；用篡改、跨会话替换、过期和重放检查状态协议；用受控批量请求检查反采集策略及正常任务误拦。状态校验、泄漏路径与规模化采集分别验收；训练隐私另按数据与训练机制评估。

\section{安全评估和发布控制}

\subsection{漏放、误拦与基率}


将安全过滤看作二分类任务，危险请求被放行为\term{漏放}{False Negative}{}，正常请求被阻止为\term{误拦}{False Positive}{}。阈值通常在二者之间权衡。设危险请求基率为 $\pi$，危险漏放概率为 $q$，正常误拦概率为 $f$，则在所有拦截中危险请求所占比例为
\begin{equation}
 \Pr(\mathrm{danger}\mid\mathrm{blocked})
 =\frac{\pi(1-q)}{\pi(1-q)+(1-\pi)f}.
\end{equation}
\begin{example}[低基率下的拦截精确率]
设请求总体含10000条请求，其中100条危险；过滤器拦截90条危险请求、误拦198条正常请求。危险召回率为90\%，正常误拦率为2\%，但288条拦截中仅90条危险，精确率为31.25\%。低危险基率使误拦的正常请求在审核队列中占较大比例。
\end{example}

若放行危险请求的代价为 $C_h$、误拦代价为 $C_b$，且有校准的后验危险概率 $p$，放行与拦截的期望成本分别为 $pC_h$ 和 $(1-p)C_b$。拦截条件为 $p>\frac{C_b}{C_h+C_b}$。这只是两动作、固定成本模型；涉及不可接受后果、人工复核和部分能力限制时，需要更完整决策规则。分类器原始分数需要先校准才能作为 $p$ 代入。

\subsection{加权指标与零失败的置信界}
对 $N$ 个样本，给定允许行为集合 $\mathcal A_i$ 与非负权重 $w_i$，可定义
\begin{equation}
 \widehat v=\frac{\sum_iw_i\mathbf1[\widehat y_i\notin\mathcal A_i]}{\sum_iw_i},
 \qquad \sum_iw_i>0.
\end{equation}
权重需预先规定，结果应同时呈现逐类失败数、分母和高严重度个案。没有正常样本的切片，其误拦率是不适用，记成零会把它读成没有误拦。还应测量要求执行控制的机会中实际执行了控制的比例；只评估过滤器本身会遗漏绕行路径。

若 $N$ 次独立同分布试验中失败数为0，真实失败率为 $p$，观测零失败的概率为 $(1-p)^N$。单侧置信度 $1-\alpha$ 的精确上界满足
\begin{equation}
 (1-p_{\mathrm{upper}})^N=\alpha
 \quad\Longrightarrow\quad
 p_{\mathrm{upper}}=1-\alpha^{\frac{1}{N}}.
\end{equation}
当 $\alpha=0.05,N=100$，上界约2.95\%；当 $N=300$，上界约0.994\%。零失败样本对应的风险上界仍为正。利用 $\conste{}^x\approx1+x$，大样本时上界近似 $-\frac{\log\alpha}{N}$，95\%置信度下约为 $\frac{3}{N}$。\footnote{这里的置信区间描述重复抽样覆盖性质。该上界适用于独立同分布样本和冻结的判定规则；自适应红队选择、相关变体与持续调阈值需要采用匹配其抽样过程的统计模型。}

\subsection{组件、链路和未知路径}
评测应从规则和参数校验，到模型单轮多轮行为，再到身份、检索、输出和工具的完整链路。\term{红队评估}{Red Teaming}{}在授权范围内主动寻找未知失效路径，结果应记录前提、真实影响和绕过的控制。仅使模型输出不合适文字与真正跨租户读取是不同后果，合并成一个攻击成功标签会掩盖这一差别。

样本需要覆盖正常易误拒请求、边界案例、不同语言与模态、长上下文以及历史事件。随机生成设置和版本应固定并报告重复次数。人工审核者也需明确判据、访问分级和暴露保护；模型评审器自身可能受注入或分布变化影响，应通过人工样本估计误差。

发布证据应绑定变化范围、风险登记、制品摘要、切片结果、权限验证和恢复能力。高严重度后果可形成独立门槛，大量正常样本的平均值掩盖不了它。受控放量应限制流量、租户、工具和资源范围；观察窗口要考虑事件基率与统计不确定性，不存在对所有场景都合适的固定百分比。

\section{监测与事件响应}

\subsection{异常调查记录}
监测需关联请求、租户、模型、策略、检索快照、工具动作和结果状态。拒答率突变可能来自攻击，也可能来自模型更新或数据变化；外部写入增加可能是业务高峰，也可能是控制绕过。告警规则应结合变更记录与分组基线；把单个指标波动直接等同于安全事件会产生大量噪声告警。

资源也是安全资产。限制输入解码大小、上下文长度、检索分支数、工具轮数、并发、输出和总截止时间，可以约束资源耗尽。取消和失败需要传播到后台任务；前端已结束而工具循环仍在运行，正是缺少这一传播的结果。对调用量异常的租户应按权限与配额处理，模型的自觉在这里没有约束力。

\subsection{遏制、保全、修复和恢复}
\term{事件响应}{Incident Response}{}是识别、遏制、调查和恢复受影响系统的组织过程。NIST SP 800-61 Rev.3（2025年）将响应准备与检测、响应、恢复放入网络安全风险管理中\cite{nist2025incidentresponse}。对模型系统，应预先具备停用工具、隔离索引、撤销凭据、切换制品与保护证据的能力；事件发生后才临时设计开关，这些能力届时都不可用。

发现疑似泄露后，先限制继续扩大的路径，例如停用外发工具和失效凭据；同时保留受控证据，包括请求时间线、来源版本、策略决策与执行账本。调查要区分已提出、已拒绝、已执行和执行结果未知的动作。没有发现成功日志不等于没有执行，工具端记录提供了另一条确认通道。

影响分析应覆盖用户、租户、数据范围和副作用，并判断缓存、记忆与派生制品是否受污染。修复位置可能在数据、模型、模板、策略或执行器。只更新提示可能掩盖一条攻击文本，错误授权原样保留；只回滚模型撤销不了已经外发的信息或已执行的业务变更。

\begin{figure}[htbp]
\centering
\begin{tikzpicture}[node distance=6mm,box/.style={draw,rounded corners,align=center,text width=29mm,minimum height=11mm,font=\small},>=stealth]
\node[box] (a) {检测与分级\\明确影响假设};
\node[box,right=of a] (b) {能力遏制\\吊销、隔离、限流};
\node[box,right=of b] (c) {证据与影响调查\\版本、执行账本};
\node[box,below=of c] (d) {根因修复\\数据、模型、策略};
\node[box,left=of d] (e) {受控恢复\\回归与范围限制};
\node[box,left=of e] (f) {复盘与治理更新\\责任、门槛、监测};
\draw[->] (a)--(b);\draw[->] (b)--(c);\draw[->] (c)--(d);\draw[->] (d)--(e);\draw[->] (e)--(f);\draw[->] (f)--(a);
\end{tikzpicture}
\caption{事件处理闭环。证据保全与遏制可并行开展，恢复以修复证据和受限能力为前提。}
\label{fig:security-response}
\end{figure}

恢复前应验证受影响控制、历史失败路径和相邻变体，并检查凭据、缓存和索引的清理结果。恢复可以分能力进行，例如先开放只读检索，继续关闭外部写入；聊天答案恢复正常不构成恢复全部工具的依据。通知、保留与披露要求按事件和组织适用规则执行，记录由谁作出相应判断。

\subsection{构造事件的时间线}


\begin{example}[内部报告外发事件的响应时序]
回到内部报告外发的构造场景。记录显示某请求在10:00提出外发建议，策略错误允许，工具在10:01执行成功，10:03出现告警。响应人员在10:05停用发送能力并撤销服务凭据，随后固定相关文档版本、动作摘要与收件范围。检测延迟从已确认执行到告警为2分钟，遏制延迟从告警到停用为2分钟；指标必须明确起止定义。

调查确认根因是“合法结构即允许”的错误策略，提示词层面只是它暴露出来的位置。修复需加入任务范围、目的地址和数据级别校验，缩小工具身份权限，并处理可能保存敏感摘要的缓存。权重回滚收不回已经送达的报告，需进入实际数据事件处置。随后以相同动作参数、不同表达、工具超时重试和缓存读取构造回归条件，分阶段恢复能力。该时间线分别记录检测、遏制、调查与恢复阶段。
\end{example}

% BEGIN GENERATED INTERVIEWS

% END GENERATED INTERVIEWS
\chapterreferences
\myendchapterdocument
